% Part: second-order-logic
% Chapter: syntax-and-semantics
% Section: language-of-sol

\documentclass[../../../include/open-logic-section]{subfiles}

\begin{document}

\olfileid{sol}{syn}{lan}

\olsection{زبان منطق مرتبهٔ دوم}

همانند منطق مرتبهٔ اول، عبارت‌های منطق مرتبهٔ دوم از واژگان پایه‌ای ساخته
می‌شوند که \emph{!!{variable}s}، \emph{!!{constant}s}،
\emph{!!{predicate}s} و گاه \emph{!!{function}s} را در بر می‌گیرد. از
آن‌ها، همراه با رابط‌های منطقی، سورها و نشانه‌های سجاوندی مانند پرانتز و
ویرگول، \emph{ترم‌ها} و~\emph{!!{formula}s} ساخته می‌شوند. تفاوت در این
است که منطق مرتبهٔ دوم، افزون بر متغیرهای اشیا، متغیرهایی برای رابطه‌ها و
تابع‌ها نیز دارد و کمیت‌گذاری بر آن‌ها را مجاز می‌داند. پس نمادهای منطقیِ
منطق مرتبهٔ دوم همان نمادهای منطق مرتبهٔ اول‌اند، به‌اضافهٔ موارد زیر:

\begin{enumerate}
\item مجموعه‌ای~!!{denumerable} از !!{variable}s رابطه‌ای مرتبهٔ دوم
  با هر مرتبهٔ~$n$: $\Obj V_0^n$، $\Obj V_1^n$، $\Obj V_2^n$، \dots
\item مجموعه‌ای~!!{denumerable} از !!{variable}s تابعی مرتبهٔ دوم
  با هر مرتبهٔ~$n$: $\Obj u_0^n$، $\Obj u_1^n$، $\Obj u_2^n$، \dots
\end{enumerate}

در منطق مرتبهٔ اول، معمولاً~!!{identity}~$\eq$ گنجانده می‌شود. نمادهای
غیرمنطقیِ زبانی چون~$\Lang{L}$ در منطق مرتبهٔ اول برای بیان هر مطلب
جالبی نقشی اساسی دارند. البته~!!{sentence}sی هستند که هیچ نماد غیرمنطقی
به کار نمی‌برند، اما تنها با~$\eq$ دشوار است مطلب جالبی گفت. در منطق
مرتبهٔ دوم، چون ذخیره‌ای نامحدود از متغیرهای رابطه‌ای و تابعی داریم، حتی
بدون ذخیره‌ای ویژه از نمادهای غیرمنطقی می‌توانیم هر آنچه را در یک زبان
مرتبهٔ اول بیان می‌کنیم، بیان کنیم.

\end{document}
